\section{การออกแบบและพัฒนาระบบ}
\subsection{สถาปัตยกรรมของระบบ}
ส่วนติดต่อผู้ใช้พัฒนาด้วย Next.js และ React ส่วนบริการประมวลผลใช้ NestJS และเข้าถึงฐานข้อมูลผ่าน Prisma ตามแฟ้มกำหนดแพ็กเกจและสคีมาในโครงการ ระบบเชื่อมบริการยืนยันตัวตน บุคลากร และข้อมูลประกอบใบลาจากภายนอก โดยการตั้งค่าปลายทางและข้อมูลลับอยู่ฝั่งเซิร์ฟเวอร์ แผนภาพสรุปองค์ประกอบหลักอยู่ในรูปที่~\ref{fig:architecture} ลูกศรสองทางหมายถึงการเรียกบริการและรับผล ส่วนลูกศรทางเดียวหมายถึงการส่งข้อมูลเข้าสู่กระบวนการรายงาน โดยบริการรายงานตรวจขอบเขตผู้ใช้กับต้นทางก่อนส่งผลให้หน้าเว็บ
\begin{figure}[ht]
\centering
\begin{tikzpicture}[
  x=1mm,y=1mm,
  box/.style={draw,rectangle,align=center,text width=32mm,minimum height=10mm,inner sep=1.5mm,font=\fontsize{12bp}{16bp}\selectfont},
  flow/.style={-{Stealth[length=1.7mm]},line width=0.5pt},
  exchange/.style={{Stealth[length=1.7mm]}-{Stealth[length=1.7mm]},line width=0.5pt}
]
\node[box] (web) at (0,0) {ผู้ใช้ผ่านเว็บระบบใบลา};
\node[box] (api) at (0,-18) {บริการระบบใบลา};
\node[box] (support) at (43,-18) {ระบบข้อมูลสนับสนุน\\วิชา ผู้สอน ที่ปรึกษา};
\node[box] (db) at (0,-39) {ฐานข้อมูลใบลา\\และ snapshot};
\node[box] (etl) at (43,-39) {ETL\\ตรวจและแปลงข้อมูล};
\node[box] (report) at (0,-60) {บริการรายงาน\\ตรวจสิทธิ์ก่อนส่งผล};
\node[box] (dw) at (43,-60) {คลังข้อมูลรายงาน};
\draw[exchange] (web) -- (api);
\draw[exchange] (api) -- (support);
\draw[exchange] (api) -- (db);
\draw[flow] (db) -- (etl);
\draw[flow] (support) -- (etl);
\draw[flow] (etl) -- (dw);
\draw[exchange] (report) -- (dw);
\draw[exchange] (report.west) -- (-20,-60) -- (-20,0) -- (web.west);
\end{tikzpicture}
\caption{องค์ประกอบและการเชื่อมต่อข้อมูลของระบบ}\label{fig:architecture}
\end{figure}
\subsubsection{เทคโนโลยีส่วนติดต่อผู้ใช้}
เว็บระบบใบลาใช้ Next.js 16.1.6 และ React 19.2.3 ส่วนเว็บระบบข้อมูลสนับสนุนกำหนด Next.js 16.2.6 และ React 19.2.4 ใน package.json ทั้งสองโครงการใช้ TypeScript สำหรับกำหนดชนิดข้อมูลของหน้าเว็บและบริการ ตัวเลขเหล่านี้เป็นรุ่นที่ประกาศในโครงการที่ตรวจ ไม่ใช่การยืนยันรุ่นของระบบที่ติดตั้งจริง \cite{repo}
ส่วนติดต่อใช้ Tailwind CSS 4 ร่วมกับคอมโพเนนต์แนวทาง shadcn/ui และ Radix UI เพื่อประกอบแบบฟอร์ม ตาราง และกล่องโต้ตอบ การแยกคอมโพเนนต์ช่วยให้หน้าสำหรับนิสิต อาจารย์ และเจ้าหน้าที่ใช้รูปแบบร่วมกันได้ ขณะที่ตรรกะเฉพาะบทบาทยังแยกในหน้าหรือบริการที่เกี่ยวข้อง
แบบฟอร์มยื่นใบลาใช้ React Hook Form จัดการค่าช่องกรอกและสถานะฟอร์ม ร่วมกับ Zod และ resolver สำหรับตรวจรูปแบบข้อมูลก่อนส่ง การตรวจฝั่งหน้าเว็บช่วยให้แจ้งข้อผิดพลาดใกล้ช่องกรอก แต่ไม่แทนการตรวจเงื่อนไขและสิทธิ์ฝั่งเซิร์ฟเวอร์ Axios ใช้รับส่งข้อมูลกับบริการ ส่วน Recharts ใช้สร้างกราฟรายงาน และส่วนเอกสารใช้ html-to-image ร่วมกับ jsPDF ในขั้นตอนส่งออกเอกสาร
\subsubsection{เทคโนโลยีบริการและการจัดเก็บ}
บริการทั้งสองโครงการใช้ NestJS 11 บน Node.js โดยแยก controller สำหรับรับคำขอ service สำหรับกฎการทำงาน และ module สำหรับประกอบความสามารถแต่ละด้าน เช่น ใบลา ข้อมูลสนับสนุน และรายงาน โครงสร้างนี้ทำให้ตรรกะการเปลี่ยนสถานะไม่จำเป็นต้องอยู่ในหน้าเว็บ และใช้ซ้ำได้จากผู้เรียกที่ได้รับสิทธิ์
Prisma 7.8 เป็นเครื่องมือเข้าถึงฐานข้อมูลที่ประกาศในทั้งสองโครงการ โดยใช้ adapter สำหรับ MariaDB driver กับฐานข้อมูลตระกูล MySQL สคีมากำหนดคีย์ ความสัมพันธ์ และข้อจำกัดข้อมูล ส่วนการเขียนคำขอกับขั้นอนุมัติใช้ธุรกรรมเพื่อรักษาความสอดคล้อง Multer รองรับการรับแฟ้มผ่านบริการฝั่งเซิร์ฟเวอร์ โดยการตรวจเจ้าของและสถานะคำขอเป็นตรรกะอีกส่วนหนึ่งของระบบใบลา
ตารางที่~\ref{tab:tech-stack} สรุปเทคโนโลยีตามหน้าที่ เพื่อแยกเฟรมเวิร์กส่วนติดต่อ เครื่องมือจัดเก็บ และเครื่องมือวิเคราะห์ออกจากกัน รายการแพ็กเกจที่มีอยู่เพียงอย่างเดียวไม่ใช่หลักฐานว่าทุกความสามารถถูกเปิดใช้ในการติดตั้งจริง
\begin{table}[ht]
\caption{เทคโนโลยีหลักและหน้าที่ในระบบ}\label{tab:tech-stack}
\centering
\begin{tabularx}{\columnwidth}{|p{30mm}|X|}
\hline
เทคโนโลยี & หน้าที่ \\ \hline
Next.js / React & หน้าเว็บและคอมโพเนนต์ผู้ใช้ \\ \hline
Tailwind / shadcn/ui & รูปแบบและองค์ประกอบหน้าจอ \\ \hline
React Hook Form / Zod & จัดการและตรวจข้อมูลแบบฟอร์ม \\ \hline
NestJS / TypeScript & บริการและกฎการทำงาน \\ \hline
Prisma / MySQL & เชื่อมข้อมูลและจัดเก็บธุรกรรม \\ \hline
Recharts & กราฟสถิติการลา \\ \hline
xlsx & อ่านข้อมูล Excel ในระบบสนับสนุน \\ \hline
Node.js / SQL & ประมวลผล ETL และคลังข้อมูล \\ \hline
\end{tabularx}
\end{table}
\subsubsection{เทคโนโลยีข้อมูลสนับสนุนและคลังข้อมูล}
บริการนำเข้าของระบบข้อมูลสนับสนุนใช้ไลบรารี xlsx อ่าน workbook และแปลงแถวเป็นข้อมูลวิชา ตอนเรียน ผู้สอน หรือที่ปรึกษา ก่อนเขียนด้วย Prisma การจัดการรูปแบบ Excel จึงเป็นหน้าที่ของบริการนำเข้า ส่วนกฎรายการซ้ำและความสัมพันธ์อาศัยทั้งตรรกะโปรแกรมและคีย์ของฐานข้อมูล
คลังข้อมูล v2 ใช้สคริปต์ JavaScript แบบ ES modules บน Node.js สำหรับ Extract--Transform--Load และใช้ mariadb driver รัน SQL โดยตรง ตารางวิเคราะห์ dw\_* กำหนดด้วย SQL แยกจากสคีมา Prisma ของธุรกรรม แม้ในโครงการยังมีสคีมาคลังข้อมูลเดิมอยู่ จึงไม่ควรอธิบายว่ารายงานปัจจุบันอ่านตารางเดิมทั้งหมด
เอกสาร ETL ระบุฐานเป้าหมาย MySQL 8.0.16 ขึ้นไปเพื่อรองรับข้อจำกัดและคำสั่งที่ใช้ การมีไดรเวอร์ชื่อ MariaDB ไม่ได้ยืนยันว่าคลังข้อมูลรองรับ MariaDB ทุกเวอร์ชัน การติดตั้งต้องตรวจชนิดและรุ่นฐานข้อมูลจริงให้ตรงกับ SQL ที่ใช้
\subsubsection{เครื่องมือพัฒนาและการติดตั้ง}
โครงการจัดเป็น pnpm monorepo แยก apps/web และ apps/api และมี lockfile ควบคุมการติดตั้ง dependency ระบบใบลาระบุ Node.js 22 และ pnpm 10.0.0 ในคู่มือและ Dockerfile โดยมี image สำหรับเว็บและบริการแยกกัน การแยกส่วนนี้ทำให้กำหนดตัวแปรแวดล้อมและพื้นที่เก็บแฟ้มของแต่ละบริการได้ตามหน้าที่
ESLint และ TypeScript ใช้ตรวจซอร์สโค้ด ส่วนชุดทดสอบสิทธิ์และ ETL ของระบบใบลาใช้ตัวทดสอบของ Node.js สำหรับระบบข้อมูลสนับสนุนมีการกำหนด Jest และ Supertest ในโครงการ การกล่าวถึงเครื่องมือเหล่านี้เป็นคำอธิบายสิ่งที่มีในโค้ด ไม่ใช่รายงานว่ารันผ่านทุกชุดในการจัดทำบทความครั้งนี้
\subsection{ระบบข้อมูลสนับสนุน}
ระบบข้อมูลสนับสนุนการลานิสิต คณะวิศวกรรมศาสตร์ กำแพงแสน (Student Leave Reference Data Management System) เรียกสั้นว่า ``ระบบข้อมูลสนับสนุน'' พัฒนาภายใต้โครงการ leave-request-components เป็นระบบจัดการข้อมูลอ้างอิงแยกจากฐานข้อมูลคำขอลา มีหน้าเว็บสำหรับจัดการรายวิชา ตอนเรียน ผู้สอน และความสัมพันธ์ที่ปรึกษา--นิสิต รวมทั้งนำเข้าข้อมูลจาก Excel บริการ NestJS เปิดให้ระบบใบลาค้นหารายวิชา เลือกตอนเรียนและผู้สอน ค้นที่ปรึกษาของนิสิต และอ่านนิสิตในความดูแลหรือรายวิชาที่สอน \cite{components}
สคีมาประกอบด้วย subjects, subject\_sections, subject\_section\_lecturers และ advisor\_student รายวิชาหนึ่งมีหลายตอนเรียนและแต่ละตอนมีผู้สอนได้หลายคน ส่วน student\_code เป็นคีย์ไม่ซ้ำในความสัมพันธ์ที่ปรึกษา บริการส่งรหัส BigInt เป็นข้อความเพื่อคงความแม่นยำของรหัสระหว่างระบบ
การนำเข้า Excel แยกชุดรายวิชากับชุดที่ปรึกษา โดยอ่านรหัสวิชา ชื่อวิชา ตอนเรียน ตารางสอน และรหัสผู้สอน หรือรหัสนิสิต ที่ปรึกษา และภาควิชา มีทั้งแบบแทนที่ชุดเดิมและแบบเพิ่มรายการใหม่ซึ่งข้ามรายการซ้ำตามคีย์ การเพิ่มข้อมูลจึงไม่เท่ากับการแก้ความสัมพันธ์เดิม ต้องเลือกกระบวนการให้ตรงกับรอบปรับปรุงข้อมูล
ระบบใบลาเชื่อม ระบบข้อมูลสนับสนุน ผ่านค่าปลายทางบริการฝั่งเซิร์ฟเวอร์ ข้อมูลนี้ใช้ทั้งประกอบแบบฟอร์มและตรวจขอบเขตปัจจุบันของที่ปรึกษาและผู้สอนก่อนอ่านรายงาน ขณะที่ข้อมูล ณ เวลายื่นเก็บใน snapshot เพื่อไม่ใช้ความสัมพันธ์ปัจจุบันแทนประวัติทั้งหมด ทั้งนี้ ระบบข้อมูลสนับสนุน ยังไม่มีมิติภาคการศึกษาและหลักสูตรในสคีมาดังกล่าว จึงต้องใช้แหล่งอ้างอิงเพิ่มเติมเมื่อต้องการประวัติรายเทอมหรือเจ้าของรายวิชาทางการ
\subsection{การยื่นคำขอและการจัดเก็บข้อมูล}
นิสิตระบุประเภทลา วันที่เริ่มและสิ้นสุด เหตุผล และรายวิชาที่ได้รับผลกระทบ พร้อมแนบหลักฐานตามเงื่อนไขของระบบ ข้อมูลคำขอเชื่อมกับนิสิต ประเภทลา และภาคการศึกษา ส่วนข้อมูลรายละเอียดประกอบเก็บผ่านแฟ้ม snapshot ที่คำขออ้างถึง ตารางหลักประกอบด้วยข้อมูลบัญชี นิสิต ใบลา ขั้นอนุมัติ และการรับทราบของผู้สอน \cite{schema}
ระบบรองรับการลาป่วยและลากิจ โดยเมื่อสร้างคำขอ เซิร์ฟเวอร์กำหนดสถานะเป็น PENDING และสร้างรายการผู้อนุมัติตามเส้นทางที่ระบบคำนวณ ขั้นแรกเป็น CURRENT ส่วนขั้นถัดไปเป็น PENDING โดยเพิ่มขั้นรองคณบดีเฉพาะเมื่อจำนวนวันลามากกว่า 15 วัน การสร้างคำขอและขั้นอนุมัติใช้ธุรกรรมเดียวกัน เพื่อรักษาความสัมพันธ์ของข้อมูล ไม่รับสถานะอนุมัติสำเร็จจากผู้ยื่นคำขอโดยตรง
\subsection{การพิจารณาและการรับทราบ}
ผู้อนุมัติพิจารณาได้เฉพาะขั้น CURRENT ที่ตรงกับตัวตนและบทบาทปัจจุบัน เมื่ออนุมัติ ระบบเปิดขั้นถัดไป หากไม่มีขั้นเหลือจึงเปลี่ยนคำขอเป็น APPROVED เมื่อปฏิเสธ คำขอเปลี่ยนเป็น REJECTED การเปลี่ยนขั้นตรวจเงื่อนไขสถานะภายในธุรกรรม เพื่อลดโอกาสที่การตัดสินพร้อมกันจะเขียนทับผลเดิม
ระบบมีส่วนอนุมัติอัตโนมัติตามกฎที่กำหนดและบันทึกแหล่งที่มาของคำตัดสินแยกจากการตัดสินด้วยตนเอง ส่วนอาจารย์ผู้สอนรับทราบใบลาที่อนุมัติแล้วผ่านความสัมพันธ์ของรายวิชาและตอนเรียน การรับทราบนี้เป็นคนละกิจกรรมกับขั้นอนุมัติและไม่ควรนำไปนับเป็นจำนวนผู้อนุมัติ
\subsection{แบบจำลองข้อมูลธุรกรรม}
ตรวจสคีมาจริงเมื่อ 18 ต.ค. 2569 ฐานใบลาและคลังข้อมูลบน Railway มีฐานละ 19 ตาราง โดยคลังข้อมูลใช้ dw\_* ส่วนระบบข้อมูลสนับสนุนบน MySQL ใน Docker มี 5 ตาราง รวมตาราง migration แผนภาพทั้งสามฐานและ SQL สำหรับสร้าง diagram อยู่ในภาคผนวกท้ายเอกสาร เส้นเชื่อมแสดงเฉพาะ foreign key ภายในแต่ละฐาน ไม่สร้างความสัมพันธ์ข้ามระบบจากชื่อรหัสที่คล้ายกัน

ฐานข้อมูลธุรกรรมแยกบัญชีผู้ใช้กับข้อมูลนิสิตเพื่อเชื่อมตัวตนกับเจ้าของคำขอ ใบลาอ้างถึงนิสิต ประเภทลา และภาคการศึกษา ขณะที่รายการอนุมัติอ้างถึงทั้งใบลาและผู้พิจารณา ความสัมพันธ์แบบหนึ่งใบลาต่อหลายขั้นช่วยเก็บลำดับ สถานะ เหตุผล และเวลาของแต่ละคำตัดสินได้โดยไม่รวมทุกขั้นไว้ในช่องข้อมูลเดียว ตารางที่~\ref{tab:operational-model} สรุปกลุ่มข้อมูลที่ใช้ในกระบวนการ
\begin{table}[ht]
\caption{กลุ่มข้อมูลหลักของฐานธุรกรรม}\label{tab:operational-model}
\centering
\begin{tabularx}{\columnwidth}{|p{23mm}|X|}
\hline
กลุ่มข้อมูล & บทบาทในกระบวนการ \\ \hline
บัญชีและนิสิต & เชื่อมตัวตนผู้ยื่นกับคำขอและประวัติ \\ \hline
ใบลา & ประเภท ช่วงวันที่ จำนวนวัน สถานะ และตำแหน่งแฟ้ม \\ \hline
ขั้นพิจารณา & ผู้รับผิดชอบ บทบาท ลำดับ ผลตัดสิน และเวลา \\ \hline
การรับทราบ & เชื่อมใบลา วิชา ตอนเรียน และผู้สอน \\ \hline
ภาคการศึกษา & กำหนดบริบทปีและเทอมของคำขอ \\ \hline
\end{tabularx}
\end{table}
การแยกข้อมูลที่ใช้ค้นหาและเชื่อมความสัมพันธ์ไว้ในฐานข้อมูล กับรายละเอียดประกอบใน snapshot ทำให้ระบบติดตามคำขอผ่านคีย์ถาวรได้ อย่างไรก็ตาม การสำรองเฉพาะฐานข้อมูลยังไม่เพียงพอหากไม่มีแฟ้มที่ถูกอ้างถึง การดูแลระบบจึงต้องครอบคลุมทั้งฐานข้อมูลและพื้นที่เก็บเอกสาร เพื่อให้เปิดรายละเอียดใบลาเดิมและตรวจหลักฐานย้อนหลังได้
\subsection{สถานะคำขอและเงื่อนไขการเปลี่ยนขั้น}
สถานะระดับคำขอตอบว่ากระบวนการเสร็จสิ้นหรือยัง ส่วนสถานะระดับขั้นตอบว่าผู้ใดมีหน้าที่พิจารณาในขณะนั้น คำขอที่ยังดำเนินการใช้ PENDING แม้บางขั้นอนุมัติแล้ว จึงไม่ควรตีความว่าคำขอทั้งใบอนุมัติเพียงเพราะพบคำตัดสิน APPROVED ในขั้นใดขั้นหนึ่ง
ตารางที่~\ref{tab:state-transition} อธิบายความสัมพันธ์ของเหตุการณ์กับสถานะ โดยเป็นคำอธิบายตรรกะจากบริการจัดการใบลา ไม่ใช่ผลทดลองการส่งคำขอจริง ผู้พิจารณาต้องผ่านเงื่อนไขตัวตน บทบาท ผู้รับผิดชอบ และสถานะพร้อมกันก่อนเปลี่ยนขั้น
\begin{table}[ht]
\caption{เหตุการณ์และการเปลี่ยนสถานะคำขอ}\label{tab:state-transition}
\centering
\begin{tabularx}{\columnwidth}{|p{25mm}|X|}
\hline
เหตุการณ์ & ผลต่อกระบวนการ \\ \hline
สร้างคำขอ & ใบลา PENDING; ขั้นแรก CURRENT \\ \hline
อนุมัติระหว่างทาง & ขั้นเดิม APPROVED; เปิดขั้นถัดไปเป็น CURRENT \\ \hline
อนุมัติขั้นสุดท้าย & คำขอเปลี่ยนเป็น APPROVED \\ \hline
ปฏิเสธขั้นปัจจุบัน & คำขอเปลี่ยนเป็น REJECTED \\ \hline
ส่งผลตัดสินซ้ำ & ต้องไม่เขียนทับขั้นที่ตัดสินเสร็จแล้ว \\ \hline
\end{tabularx}
\end{table}
การอนุมัติด้วยตนเองและกฎอัตโนมัติอาจเข้าถึงขั้นเดียวกันในเวลาใกล้กัน การตรวจและยึดสิทธิ์เปลี่ยนขั้นภายในธุรกรรมจึงเป็นส่วนสำคัญของการรักษาคำตัดสินเพียงผลเดียว ข้อมูลแหล่งที่มาของคำตัดสินช่วยแยกว่าผลนั้นเกิดจากผู้พิจารณาหรือกฎ เพื่อให้ตรวจสอบย้อนหลังได้
\subsection{การควบคุมหลักฐานและเอกสาร}
หลักฐานการลาต้องเข้าถึงผ่านการยืนยันตัวตนและการตรวจสิทธิ์ของคำขอ เจ้าของคำขออัปโหลดหรือแทนที่หลักฐานได้ตามเงื่อนไขก่อนมีผลการพิจารณา เว็บเรียกเอกสารแนบผ่านตัวกลางฝั่งเซิร์ฟเวอร์ ส่วนลิงก์เอกสารสาธารณะไม่ส่งคืนหลักฐานส่วนตัว จึงแยกการตรวจสอบเอกสารใบลาออกจากการเผยแพร่สิ่งแนบ
\subsection{รายงานตามบทบาท}
นิสิตดูประวัติของตน ที่ปรึกษาดูนิสิตในความรับผิดชอบ ผู้สอนดูรายวิชาที่เกี่ยวข้อง ส่วนหัวหน้าภาค ผู้บริหาร และเจ้าหน้าที่ดูข้อมูลตามขอบเขตที่ได้รับสิทธิ์ หน้ารายงานมีการกรองปีการศึกษา ภาคการศึกษา และมิติที่เกี่ยวข้องกับบทบาท
คลังข้อมูลใช้ข้อมูลระดับคำขอสำหรับจำนวนใบลาและวันลา และข้อมูลระดับรายวิชาสำหรับสถิติวิชา เพื่อหลีกเลี่ยงการรวมวันลาซ้ำ การโหลดเป็น snapshot แบบสั่งรันและมีเวลาการโหลดสำเร็จล่าสุด ผู้ใช้จึงต้องพิจารณาความสดของข้อมูลก่อนใช้ประกอบการตัดสินใจ

\subsection{โครงสร้างคลังข้อมูลและนิยามตัวชี้วัด}
คลังข้อมูลรุ่น v2 ใช้ตารางตระกูล dw\_* แยกจากโครงสร้าง wh\_* เดิม โดยมีตารางข้อเท็จจริงหลักสามระดับ ได้แก่ หนึ่งใบลาต่อแถวใน dw\_fact\_leave\_request หนึ่งใบลา--รายวิชาที่เปิดสอนต่อแถวใน dw\_fact\_leave\_subject และหนึ่งขั้นพิจารณาต่อแถวใน dw\_fact\_approval\_step มิตินิสิต บุคลากร ประเภทลา ภาคการศึกษา หน่วยงาน หลักสูตร วิชา และตอนเรียนเชื่อมกับข้อเท็จจริง ส่วนผู้สอนร่วมเชื่อมผ่านตารางความสัมพันธ์หลายต่อหลาย
จำนวนใบลานับจากรหัสคำขอไม่ซ้ำและรวมวันลาที่ระดับคำขอ ตัวอย่างเช่น ใบลาหนึ่งใบจำนวน 3 วันกระทบ 2 วิชา ยังคงนับเป็น 1 ใบและ 3 วัน แม้มีข้อเท็จจริงรายวิชา 2 แถว กรณีมีใบลาสองใบจำนวนวันเท่ากัน ต้องรวมทั้งสองใบ ไม่ตัดซ้ำด้วยค่าจำนวนวัน การตัดซ้ำใช้รหัสคำขอก่อนรวมยอด วันลาในปฏิทินแยกจากวันขาดเรียนจริงที่ต้องตรวจวันที่และตารางสอนก่อนรวม ส่วนเวลาพิจารณาคำนวณจากเวลาที่ขั้นพร้อมให้ตัดสินถึงเวลาตัดสิน ไม่ใช้เวลาสร้างรายการแทนโดยอัตโนมัติ
รายงานค่าเฉลี่ยรายภาคการศึกษาต้องรวมเทอมที่ไม่มีการลาเป็นศูนย์ในตัวหาร การจัดอันดับคงตัวเลือกปีและภาคการศึกษา ข้อมูลสำหรับ box plot ของผู้สอนใช้ยอดรายวิชาต่อเทอม ไม่ใช้ระยะเวลาของใบลาแต่ละใบ การกำหนดหน่วยวิเคราะห์เหล่านี้ทำให้กราฟและตารางตอบคำถามเดียวกันได้อย่างสอดคล้อง
\subsection{กระบวนการ ETL และความสดของรายงาน}
ขั้น Extract อ่านฐานธุรกรรมและ JSON snapshot พร้อมข้อมูลอ้างอิงจากบริการบุคลากรและ ระบบข้อมูลสนับสนุน ขั้น Transform แปลงรหัส เชื่อมนิสิต ผู้สอน ที่ปรึกษา วิชา และภาคการศึกษา รวมทั้งตรวจค่าที่ขาดหรือกำกวม ขั้น Load ปรับข้อมูลด้วยคีย์ต้นทางเดิมและแทนชุดข้อมูลลูกของคำขอภายในธุรกรรมเดียว การรันซ้ำจึงคงคีย์และไม่เพิ่มใบลาเดิมซ้ำ ส่วน reconcile ใช้จัดการรายการที่หายจาก snapshot ครบชุด
กระบวนการปัจจุบันเป็น full snapshot ในหน่วยความจำ ไม่ใช่การดึงเฉพาะส่วนที่เปลี่ยนแบบ CDC มีล็อกผู้เขียนข้อมูลครั้งละหนึ่งราย บันทึก batch และตรวจคุณภาพก่อน commit หากโหลดล้มเหลวจะ rollback ชุดข้อมูลและบันทึกความล้มเหลว ระบบแสดงเวลาของ batch ที่สำเร็จล่าสุดเพื่อให้ผู้อ่านประเมินความสดได้
บริการรายงานปัจจุบันอ่าน dw\_* และตรวจตัวตน บทบาท ตลอดจนรายชื่อนิสิตหรือตอนเรียนปัจจุบันจากต้นทางก่อนจำกัดข้อมูล การมีข้อมูลใน warehouse จึงไม่ใช่หลักฐานว่าผู้ใช้มีสิทธิ์อ่าน และข้อมูลอ้างอิงที่เก็บไว้เพื่อวิเคราะห์ย้อนหลังไม่แทนการตรวจสิทธิ์ปัจจุบัน
